\section{结语：统一称号不能抹平多政权遗产}

元把大都、江南、边疆与海域置于一个帝国框架，并未消除这些地区在资源、文书能力和地方社会上的差异。行省、纸钞、漕运、法律与宗教管理能够扩大调度范围，也可能在征收、折纳、交通中断和财政信用受损时把压力转给家庭与地方社群。元末的危机应由这些不均衡关系、河工、战乱与地方政治的交错来理解，而非归于某一制度或群体的单独失败。

回看全卷，宋、辽、西夏、金、南宋、蒙古与元始终在彼此塑造的道路、市场、边界和迁徙网络中行动。政权更替改变了这些网络的主导者，却没有使它们的代价、地方差异或史料沉默消失。随后的跨政权综合将据此比较各国如何获得治理能力，又如何在竞争、征服与社会承受力面前显出限度。
